% !TeX root = ../main.tex

\chapter{排名母体与组合构成机制}
\section{排名变化的描述性幅度}
\begin{table}[htbp]
\centering
\caption{成对调整的描述性变化}
\renewcommand{\arraystretch}{1.18}
\small
\setlength{\tabcolsep}{4.5pt}
\begin{tabularx}{\textwidth}{@{}p{1.45cm}*{4}{C}@{}}
\toprule
对照 & 年化变化/百分点 & 波动变化/百分点 & 回撤变化/百分点 & 夏普变化 \\
\midrule
E$\to$P1 & 3.48 & -2.43 & 1.24 & 0.140 \\
P1$\to$P2 & -2.26 & 0.95 & 4.72 & -0.080 \\
P1$\to$P3 & -6.46 & -8.53 & 36.13 & 0.063 \\
P2$\to$P4 & -6.53 & -8.72 & 29.29 & 0.026 \\
\bottomrule
\end{tabularx}
\end{table}

表10中的E到P1是整套组合规则变化，P1到P2才对应固定历史池和得分后的排名母体调整；P1到P3、P2到P4对应同一底层组合的仓位管理。回撤变化取“新值$-$旧值”，正值表示负回撤收窄，不能把它误读为风险增加。收益、波动和回撤变化用百分点，不是相对百分比。

\begin{figure}[htbp]
\centering
\includegraphics[width=\linewidth]{fig10}
\caption{收益—波动位置与成对风控调整}
\par\smallskip
\begin{minipage}{0.97\linewidth}\small 注：箭头分别连接P1到P3、P2到P4；该图不表示均值—方差有效前沿。资料来源：CSMAR，作者计算。\end{minipage}
\end{figure}

P1与P2的最终名单在成熟期74个调仓日均发生变化。两种排名既改变核心范围，也改变扩展候选，最后通过统一三因子截取改变持仓。全截面排名不能独立于最终组合规则解释，其收益差是相同程序其余环节下的实施效果。

\begin{figure}[htbp]
\centering
\includegraphics[width=\linewidth]{fig11}
\caption{排名调整的相对净值与持仓集合相似度}
\par\smallskip
\begin{minipage}{0.97\linewidth}\small 注：相对净值为P2净值除以P1净值；Jaccard为两名单交集数除以并集数。该相似度不是成交换手率。资料来源：CSMAR，作者计算。\end{minipage}
\end{figure}

\section{持股数量、缺失收益与当前成长率}
\begin{table}[htbp]
\centering
\caption{持仓覆盖与当前成长率诊断}
\renewcommand{\arraystretch}{1.18}
\small
\setlength{\tabcolsep}{4.5pt}
\begin{tabularx}{\textwidth}{@{}p{1.55cm}*{5}{C}@{}}
\toprule
方案 & 平均持股/只 & 空仓信号/次 & 记录数/条 & 收益缺失月/个 & 当期非正成长/条 \\
\midrule
E & 5.96 & 8 & 560 & 5 & 0 \\
P1 & 14.27 & 0 & 1341 & 16 & 0 \\
P2 & 14.53 & 0 & 1366 & 36 & 17 \\
\bottomrule
\end{tabularx}
\end{table}

E平均持股5.96只、8个空仓信号；P1、P2分别为14.27和14.53只，均无空仓，但这不等于每季固定15只，初始期仍可能不足。E成熟期交集全部纳入，最大持股18只；P1、P2成熟期由统一得分截于15只。

目标持股多于有效月收益的月份，E有5个，P1有16个，P2有36个。所有路径沿用相同缺失收益再平均方式，但扩展候选改变了缺失事件发生频率，因而不能认为缺失收益对路径比较必然无影响。未有收益记录的股票不应被解释为已经以合理价格卖出；停退市真实损失尚未还原。

P2有17条股票—季度记录的当前成长率非正。这与20季至少19季为正的历史条件并不矛盾，也与正PE条件不同。全市场前50\%在盈利普遍下降时可以包含负成长企业，因此本路径保留的是较宽的历史持续性约束，而非要求每一只持仓当季都在增长。P1本次未出现这类持仓，亦不能据此把其程序门槛改写为“当前成长必须为正”。

\section{候选核心与最终核心}
\begin{table}[htbp]
\centering
\caption{成熟期核心候选与最终持仓的区别}
\renewcommand{\arraystretch}{1.18}
\small
\setlength{\tabcolsep}{4.5pt}
\begin{tabularx}{\textwidth}{@{}p{1.55cm}*{4}{C}@{}}
\toprule
路径 & 候选核心/条 & 最终核心/条 & 候选核心保留率/\% & 持仓核心占比/\% \\
\midrule
P1 & 439 & 309 & 70.39 & 28.48 \\
P2 & 552 & 137 & 24.82 & 12.34 \\
\bottomrule
\end{tabularx}
\end{table}

P2候选核心由439条增至552条，但最终保留核心由309条降至137条。候选核心条数跨季度累计，不是独立公司数；候选核心保留率以候选核心为分母，最终持仓核心占比以全部最终成熟期持仓记录为分母，两个比例不可混用。

这一变化与补位池从1,097条扩大到3,310条、最终又统一按得分截取的机制一致。扩展候选可以挤出核心，较大的核心池不保证最终核心比例提高。当前正式路径保留了用户指定的组合实现，因此本文不把结果称为“在不删任何核心股票的前提下改进”。

\begin{figure}[htbp]
\centering
\includegraphics[width=\linewidth]{fig12}
\caption{目标持股与实际持仓核心比例}
\par\smallskip
\begin{minipage}{0.97\linewidth}\small 注：核心比例仅用最终持仓标记计算；成熟期从2008年第一笔成熟期信号开始，基础原规则没有补位分层。资料来源：CSMAR，作者计算。\end{minipage}
\end{figure}

\section{为什么补位与核心保留必须分开讨论}
扩展候选的经济动机是缓解严格历史条件和两项30\%交集造成的持仓过少：在仍保留持续成长、相对估值和正PE约束的前提下，从较宽范围寻找替代股票。质量、动量与规模用于统一排序，它们只是对候选之间作选择，不能替代最初的成长与估值筛选。因此，这些改进仍保留原选股思想，但已经不是原交集的简单复制。

代码目前把候选核心与扩展候选合并，再统一按得分取最多15只。只要扩展候选得分更高，就可能排在核心前面；“核心”这个标签本身没有优先保留权。P2核心候选增加而最终核心减少，正好提醒研究者：扩大一个上游集合不保证扩大下游结果，最终截取规则同样重要。若日后要研究核心优先，应独立保存并另作对照，不能只修改报告文字来宣称已经实现。

持股数量增加也不等于完成风险分散证明。多只股票可能共享行业、规模或盈利周期暴露，现有结果只直接说明名单容量和部分风险统计的变化。本文没有据此宣称行业中性或独立风险已消除；原研报的行业比例不能直接代替当前组合的行业证据。
